\section{Methodology}

\subsection{Experimental design}

We run five families of experiments:
\begin{itemize}
  \item monolingual baselines for English, Norwegian and Danish,
  \item zero-shot transfer to Danish from English and Norwegian,
  \item transfer fine-tuning, where a source-trained model is further trained on Danish,
  \item Danish-only supervision at multiple data budgets (10\%, 25\%, 50\%, 100\%), and
  \item linguistic-distance analysis based on WALS typological feature vectors.
\end{itemize}

The label space and Danish target splits stay fixed throughout. What varies are source language, encoder, and Danish supervision budget — three clean axes that keep the interpretation tractable.

\subsection{Models and checkpoints}

Two multilingual encoders serve as the backbone for token-classification:
\begin{itemize}
  \item \textbf{mBERT}: \texttt{bert-base-multilingual-cased} \citep{devlin2019bert} — 12 layers, 110M parameters, trained on 104 languages.
  \item \textbf{XLM-R}: \texttt{xlm-roberta-base} \citep{conneau2020xlmr} — 12 layers, 125M parameters, trained on 100 languages with a substantially larger pretraining corpus than mBERT.
\end{itemize}

An English-only BERT baseline (\texttt{bert-base-cased}) provides a single-language reference point. All other experiments attach a linear token-classification head for BIO tagging on top of the multilingual encoders.

For fine-tuned transfer (EN→DA FT, NO→DA FT) we initialize from the best zero-shot checkpoint and continue training on Danish. These models have therefore been exposed to labeled data from both the source and the target language.

\subsection{Training hyperparameters}

All runs share the same configuration unless otherwise stated:

\begin{table}[t]
\centering
\small
\begin{tabular}{ll}
\toprule
Hyperparameter & Value \\
\midrule
Learning rate          & $2\times10^{-5}$ (AdamW) \\
Train batch size       & 8 \\
Eval batch size        & 16 \\
Max sequence length    & 256 tokens \\
Training epochs        & 4 \\
Weight decay           & 0.01 \\
Random seed            & 42 \\
Early stopping         & Best checkpoint by dev F1 \\
\bottomrule
\end{tabular}
\caption{Shared hyperparameters used across all experiments. Multi-seed replication and its implications are discussed in Section~\ref{sec:discussion}.}
\label{tab:hyperparams}
\end{table}

We deliberately avoided per-language hyperparameter tuning, applying the same configuration to every run. That choice can leave performance on the table for individual systems, but it keeps comparisons interpretable.

\subsection{Evaluation}

The primary metric is strict span-level F1 as computed by \texttt{span\_f1.py}, requiring an exact match on both entity boundaries and type. We also report:
\begin{itemize}
  \item unlabeled span F1 (exact boundaries, any label),
  \item loose span F1 (partial overlap, matching label).
\end{itemize}
These additional views help separate boundary errors from label confusions from entirely missed spans, giving a more complete picture of where each system struggles.

Because the English test split is masked (no gold labels), test-set comparisons focus on Danish and Norwegian. Danish-target models are evaluated on the Danish DDT test split; Norwegian source baselines on the NorNE test split.

\subsection{Statistical testing}

To assess the English--Norwegian contrast we apply paired resampling tests (paired bootstrap, 10{,}000 iterations) to dev and test prediction files. With only a few seeds, these tests serve as a practical reliability check, not a substitute for broader replication. More seeds would still be preferable.
